\section{Conclusion}

We introduced persistent navigation in evolving worlds, where an embodied agent
must act while task-relevant changes can occur outside its view before or during
navigation. \evolvingworldnav{} closes a memory--belief--action loop: its
\pfdnav{} core separates persistence from relocation, retains uncertainty over
scene-specific candidates, converts visibility-qualified observations into
posterior evidence, and replans from the updated belief. \pfdbench{} grounds
this question in human-trace-driven HSSD histories and paired static, routine,
and random worlds. The strongest gains occur under learnable routine evolution,
while lower performance on matched random transitions defines the current
boundary of predictive benefit. Results on held-out scenes, other embodied
benchmarks, and physical robots further suggest that the belief-and-evidence
interface transfers beyond the controlled routine setting. Overall, navigation
in an evolving world requires not only remembering or predicting a destination,
but testing that belief against new physical evidence and revising action when
the world no longer supports it.
